\section{Theoretical Analysis}
\label{sec:theory}

\subsection{Feasible-space restriction}

\begin{proposition}
\hfill\break
For every contiguous platoon partition \(\Pi\), the platoon-constrained
optimality gap satisfies \(G(\Pi)\ge 0\).
\end{proposition}

\begin{proof}
Every sequence satisfying the platoon-continuity constraints also satisfies
the original FIFO constraints. Hence \(\pseqset\subseteq\seqset\), and
minimization over the restricted set cannot improve the objective.
\end{proof}

\subsection{A FIFO-preserving local repair}

Consider an internal platoon link
\begin{equation}
e=(a,a')=((l,i),(l,i+1)),
\end{equation}
where the FIFO index \(i\) starts from one. If the two vehicles are not
adjacent in a current FIFO-feasible sequence, that sequence has the form
\begin{equation}
s=[P,a,B,a',R],
\end{equation}
where \(B=(b_1,\ldots,b_m)\) is nonempty. Because \(a'\) is the immediate
FIFO successor of \(a\), no vehicle in \(B\) belongs to approach \(l\). The
local repair moves \(a'\) immediately behind \(a\):
\begin{equation}
s'=[P,a,a',B,R].
\label{eq:local-repair}
\end{equation}
Let
\begin{equation}
d_e=r_{a'}-r_a,
\qquad
q_e=\max\{d_e,\hF\},
\qquad
\Delta_h=\hS-\hF,
\end{equation}
and let
\begin{equation}
M_e=m+|R|+1
\end{equation}
be the number of vehicles whose passing times can change in this repair.

\begin{lemma}[Indexed local perturbation bound]
\label{lem:indexed-local-repair}
The repaired sequence \(s'\) is FIFO-feasible and its total-delay change
satisfies
\begin{equation}
J(s')-J(s)
\le
M_e\pospart{d_e-\hF}-2\Delta_h,
\label{eq:indexed-local-bound}
\end{equation}
where \(J(s)=ND(s)\). Moreover,
\begin{equation}
M_e\le N-i.
\label{eq:indexed-affected-count}
\end{equation}
\end{lemma}

\begin{proof}
Moving \(a'\) left until it immediately follows its FIFO predecessor crosses
no vehicle from approach \(l\). The relative order of every other approach is
unchanged, so \(s'\) remains FIFO-feasible.

Let \(T=C_a(s)=C_a(s')\). Since \(T\ge r_a\),
\begin{equation}
C_{a'}(s')
=\max\{r_{a'},T+\hF\}
\le T+q_e.
\end{equation}
In the original sequence, the transitions \(a\to b_1\) and
\(b_m\to a'\) are between different approaches, while every transition
inside \(B\) requires at least \(\hF\). Therefore
\begin{equation}
C_{a'}(s)
\ge T+2\hS+(m-1)\hF,
\end{equation}
and hence
\begin{equation}
C_{a'}(s')-C_{a'}(s)
\le q_e-2\hS-(m-1)\hF.
\label{eq:moved-vehicle-change}
\end{equation}

For \(b_1\), monotonicity of the maximum operator gives
\begin{equation}
C_{b_1}(s')
\le \max\{r_{b_1},T+\hS+q_e\}
\le C_{b_1}(s)+q_e.
\end{equation}
The predecessor pairs inside \(B\) are unchanged, so induction through the
completion-time recursion yields
\begin{equation}
C_{b_j}(s')-C_{b_j}(s)\le q_e,
\qquad j=1,\ldots,m.
\label{eq:block-vehicle-change}
\end{equation}

If \(R\) is nonempty, let \(x\) be its first vehicle. The predecessor of
\(x\) changes from \(a'\) to \(b_m\). Using
\(C_{b_m}(s')\le C_{b_m}(s)+q_e\),
\(C_{a'}(s)\ge C_{b_m}(s)+\hS\),
\(h(b_m,x)\le\hS\), and \(h(a',x)\ge\hF\), we obtain
\begin{equation}
C_{b_m}(s')+h(b_m,x)
\le
C_{a'}(s)+h(a',x)+(q_e-\hF).
\end{equation}
Because \(q_e-\hF=\pospart{d_e-\hF}\ge0\), another induction through the
unchanged suffix gives
\begin{equation}
C_v(s')-C_v(s)
\le q_e-\hF,
\qquad v\in R.
\label{eq:suffix-vehicle-change}
\end{equation}
Equations~\eqref{eq:moved-vehicle-change}--\eqref{eq:suffix-vehicle-change}
sum to
\begin{align}
J(s')-J(s)
&\le q_e-2\hS-(m-1)\hF+mq_e+|R|(q_e-\hF)\\
&=(m+|R|+1)(q_e-\hF)-2(\hS-\hF),
\end{align}
which is Eq.~\eqref{eq:indexed-local-bound}.

Finally, the decomposition of \(s\) gives
\begin{equation}
M_e=N-|P|-1.
\end{equation}
FIFO feasibility places the \(i-1\) predecessors of \(a=(l,i)\) inside
\(P\), so \(|P|\ge i-1\). Consequently \(M_e\le N-i\), proving
Eq.~\eqref{eq:indexed-affected-count}.
\end{proof}

\subsection{Global repair and the indexed loss bound}

For each platoon, list its vehicles in FIFO order and process its internal
links from left to right. When a processed link is not adjacent, apply
Eq.~\eqref{eq:local-repair}; otherwise do nothing. The platoons themselves
may be processed in any fixed order.

\begin{proposition}[Validity of the global repair]
\label{prop:global-repair}
Starting from any sequence in \(\seqset\), the preceding repair procedure
terminates after at most \(N-K\) local repairs and produces a sequence in
\(\pseqset\). No previously established internal platoon adjacency is
destroyed.
\end{proposition}

\begin{proof}
There are \(N-K\) internal platoon links, and each is visited once. A local
repair preserves FIFO feasibility by Lemma~\ref{lem:indexed-local-repair}.
Within a platoon, left-to-right processing ensures that when \(a'\) is moved,
the link between \(a'\) and its FIFO successor has not yet been established.
The move therefore cannot destroy a previously repaired link of that
platoon. It also cannot split an established link of another platoon: any
such link lies wholly inside \(P\), \(B\), or \(R\), whose internal orders
and adjacencies are unchanged. Thus every established adjacency persists.
After all links are visited, every platoon is a consecutive block, so the
final sequence belongs to \(\pseqset\).
\end{proof}

For an internal link \(e=((l,i_e),(l,i_e+1))\), define its FIFO-index weight
\begin{equation}
\alpha_e=\frac{N-i_e}{N}.
\label{eq:index-weight}
\end{equation}

\begin{theorem}[FIFO-indexed platoon loss bound]
\label{thm:indexed-loss-bound}
For any finite number of approaches, any positive vehicle counts, any
within-approach nondecreasing release times, and any contiguous platoon
partition \(\Pi\),
\begin{equation}
0\le G(\Pi)\le \widehat{G}(\Pi),
\label{eq:indexed-global-bound}
\end{equation}
where
\begin{equation}
\widehat{G}(\Pi)
=
\sum_{e\in\mathcal{E}_{\Pi}}
\pospart{
\alpha_e\pospart{d_e-\hF}
-\frac{2(\hS-\hF)}{N}
}.
\label{eq:indexed-partition-bound}
\end{equation}
\end{theorem}

\begin{proof}
Let \(s^0\) be an unrestricted optimal sequence. Applying
Proposition~\ref{prop:global-repair} produces the finite repair sequence
\begin{equation}
s^0,s^1,\ldots,s^T,
\qquad s^T\in\pseqset.
\end{equation}
If \(e_t\) is the link
repaired at step \(t\), Lemma~\ref{lem:indexed-local-repair} gives
\begin{equation}
J(s^t)-J(s^{t-1})
\le
M_{e_t}\pospart{d_{e_t}-\hF}-2(\hS-\hF).
\end{equation}
Telescoping, replacing each increment by its positive part, and using
\(D^*_{\Pi}\le D(s^T)\) yield
\begin{equation}
G(\Pi)
\le
\frac{1}{N}
\sum_{e\in\mathcal Q}
\pospart{
M_e\pospart{d_e-\hF}-2(\hS-\hF)
},
\end{equation}
where \(\mathcal Q\subseteq\mathcal E_\Pi\) contains the links that actually
required repair. By Eq.~\eqref{eq:indexed-affected-count} and monotonicity of
the positive-part operator,
\begin{equation}
\frac{1}{N}
\pospart{
M_e\pospart{d_e-\hF}-2(\hS-\hF)
}
\le
\pospart{
\frac{N-i_e}{N}\pospart{d_e-\hF}
-\frac{2(\hS-\hF)}{N}
}.
\end{equation}
Adding the nonnegative terms for links outside \(\mathcal Q\) proves the
upper bound. The lower bound follows from feasible-space restriction.
\end{proof}

\subsection{Rule-level optimality-gap bound}

The preceding result expresses the optimality-gap upper bound using the
release-time gap and FIFO position of each internal platoon link. We next
derive a general bound that depends directly on the two parameters used by a
rule-based platoon-formation method. Let \(\delta\) denote the critical
headway threshold and let \(P_{\max}\) denote the maximum number of vehicles
allowed in one platoon. A partition generated by this rule satisfies
\begin{equation}
d_e\le\delta,
\qquad e\in\mathcal E_\Pi,
\label{eq:rule-headway-condition}
\end{equation}
and every platoon contains at most \(P_{\max}\) vehicles.

Because \(i_e\ge1\), the FIFO-index weight satisfies
\begin{equation}
\alpha_e=\frac{N-i_e}{N}\le\frac{N-1}{N}.
\end{equation}
Together with Eq.~\eqref{eq:rule-headway-condition}, each term in
Eq.~\eqref{eq:indexed-partition-bound} is bounded by
\begin{equation}
\pospart{
\alpha_e\pospart{d_e-\hF}
-\frac{2(\hS-\hF)}{N}
}
\le
\pospart{
\frac{N-1}{N}\pospart{\delta-\hF}
-\frac{2(\hS-\hF)}{N}
}.
\label{eq:uniform-link-bound}
\end{equation}

If the partition contains \(K\) platoons, then
\(|\mathcal E_\Pi|=N-K\). Since no platoon contains more than
\(P_{\max}\) vehicles,
\begin{equation}
K\ge\left\lceil\frac{N}{P_{\max}}\right\rceil,
\end{equation}
and therefore
\begin{equation}
|\mathcal E_\Pi|
\le
N-\left\lceil\frac{N}{P_{\max}}\right\rceil.
\label{eq:rule-link-count-bound}
\end{equation}

\begin{corollary}[Rule-level platoon loss bound]
\label{cor:rule-level-loss-bound}
For any platoon partition generated using critical headway threshold
\(\delta\) and maximum platoon size \(P_{\max}\),
\begin{equation}
0\le G(\Pi)
\le \widehat{G}(\Pi)
\le \widehat{G}(\delta,P_{\max}),
\label{eq:rule-level-gap-chain}
\end{equation}
where
\begin{equation}
\widehat{G}(\delta,P_{\max})
=
\left(
N-\left\lceil\frac{N}{P_{\max}}\right\rceil
\right)
\pospart{
\frac{N-1}{N}\pospart{\delta-\hF}
-\frac{2(\hS-\hF)}{N}
}.
\label{eq:rule-level-gap-bound}
\end{equation}
\end{corollary}

\begin{proof}
Equation~\eqref{eq:uniform-link-bound} provides a common upper bound for
every internal-link contribution in \(\widehat{G}(\Pi)\), while
Eq.~\eqref{eq:rule-link-count-bound} bounds the number of such contributions.
Multiplying the two bounds and applying
Theorem~\ref{thm:indexed-loss-bound} gives
Eq.~\eqref{eq:rule-level-gap-chain}.
\end{proof}

For given traffic and safety parameters \(N\), \(\hF\), and \(\hS\),
Eq.~\eqref{eq:rule-level-gap-bound} identifies two controllable factors in
the platooning-induced optimality gap. The critical headway threshold
\(\delta\) limits the contribution associated with each internal platoon
link, while the maximum platoon size \(P_{\max}\) limits how many internal
links can accumulate. Increasing either parameter permits more aggressive
aggregation and can increase the upper bound. This result motivates a
rule-based platooning method that forms a platoon only when the consecutive
vehicle headway does not exceed \(\delta\) and the current platoon size has
not reached \(P_{\max}\). The rule retains the computational simplicity of
critical-headway platooning while explicitly controlling the second factor
identified by the optimality-gap analysis.

\begin{corollary}[Index-dependent zero-loss condition]
\label{cor:indexed-zero-loss}
If every internal platoon link
\(e=((l,i_e),(l,i_e+1))\) satisfies
\begin{equation}
d_e
\le
\hF+\frac{2(\hS-\hF)}{N-i_e},
\label{eq:indexed-zero-loss-condition}
\end{equation}
then \(G(\Pi)=0\).
\end{corollary}

\begin{proof}
Condition~\eqref{eq:indexed-zero-loss-condition} makes every summand in
Eq.~\eqref{eq:indexed-partition-bound} zero. Theorem~\ref{thm:indexed-loss-bound}
and \(G(\Pi)\ge0\) then imply \(G(\Pi)=0\).
\end{proof}

The denominator \(N-i_e\) increases the admissible zero-loss release gap for
links appearing later in an approach queue. This index dependence follows
from FIFO structure alone and does not require solving either the original or
the platoon-constrained scheduling problem.
